% BHU PYQ -- Transcribed & Corrected
% Paper: MATMJ-31: Linear Algebra
% Exam: B.A. / B.Sc. (Hons.) Semester III Examination, 2025-26 (NEP)
% Department: Mathematics

\documentclass[11pt,a4paper]{article}
\usepackage[a4paper,top=2.5cm,bottom=2.5cm,left=2.8cm,right=2.8cm,headheight=14pt]{geometry}
\usepackage{amsmath,amssymb}
\usepackage[T1]{fontenc}
\usepackage[utf8]{inputenc}
\usepackage{lmodern,microtype,fancyhdr,enumitem,hyperref,lastpage,parskip}

\hypersetup{
    colorlinks=true,
    urlcolor=black,
    linkcolor=black,
    pdftitle={MATMJ31: Linear Algebra -- BHU B.Sc. Sem III 2025-26 (NEP)}
}

\pagestyle{fancy}
\fancyhf{}
\renewcommand{\headrulewidth}{0.4pt}
\renewcommand{\footrulewidth}{0.4pt}
\fancyhead[L]{\small\textit{Banaras Hindu University}}
\fancyhead[R]{\small\textit{MATMJ-31: Linear Algebra}}
\fancyfoot[C]{\small Page \thepage\ of \pageref{LastPage}}

\newlist{parts}{enumerate}{2}
\setlist[parts,1]{label=(\alph*),leftmargin=*,itemsep=4pt,topsep=2pt}
\setlist[parts,2]{label=(\roman*),leftmargin=*,itemsep=2pt,topsep=2pt}
\newcommand{\pts}[1]{\hfill\textit{[#1]}}

\begin{document}

\begin{center}
  {\large\bfseries BANARAS HINDU UNIVERSITY}\\[2pt]
  {\normalsize B.A. / B.Sc. (Hons.) --- Semester III Examination, 2025-26 (NEP)}\\[6pt]
  \rule{0.8\linewidth}{0.5pt}\\[6pt]
  {\large\bfseries Mathematics}\\[4pt]
  {\large\bfseries Paper Code: MATMJ-31 : Linear Algebra}\\[4pt]
  {\normalsize Time: 2:30 Hours\hfill Maximum Marks: 70}
\end{center}

\rule{\linewidth}{0.4pt}

\smallskip
\noindent\textit{Instructions: Answer five questions including Question No. 1 which is compulsory. All symbols have their usual meanings.}

\bigskip

\noindent\textbf{Question 1.} Short Answer Questions (Compulsory):\pts{22}
\begin{parts}
  \item Define vector space over a field $F$. Let $V$ be a vector space over field $F$, then show that:\pts{2}
  \begin{parts}
    \item $a \cdot \mathbf{0} = \mathbf{0}$, where $a \in F$ and $\mathbf{0}$ is the zero vector.
    \item $(-a)x = -(ax) = a(-x)$, where $a \in F, x \in V$.
  \end{parts}
  \item Define sub-space of a vector space. Is the set $W = \{f(x) \in F[x] \mid 2f(0) = f(1)\} \subset F[x]$, a subspace of $F[x]$ (the vector space of all polynomials over $F$)?\pts{2}
  \item Is the set $\{(1, 1, 1), (-1, 0, 1), (0, -2, 1)\}$ of vectors in $\mathbb{R}^3$ over $\mathbb{R}$ linearly dependent?\pts{2}
  \item Let $V$ be a 3-dimensional vector space with $A$ and $B$ as its subspaces of dimensions 2 and 1 respectively. If $A \cap B = \{\mathbf{0}\}$, then show that $V = A \oplus B$.\pts{2}
  \item Show that the range space of a linear transformation $T: V \to U$ is a subspace of $U$.\pts{2}
  \item Show that the diagonal elements of a skew-Hermitian matrix are either purely imaginary or zero.\pts{2}
  \item Define Hermitian matrix. If $A$ and $B$ are Hermitian matrices, then show that $(A + B)$ is also Hermitian.\pts{2}
  \item Define invariant subspace of a vector space. Let $S = \{W_i\}_{i \in I}$ be a family of $T$-invariant subspaces of a vector space $V$, then show that $W = \bigcap_{i \in I} W_i$ is also $T$-invariant.\pts{2}
  \item Define non-singular linear transformation. Show that a linear transformation $T: V \to U$ is non-singular if and only if the image of a linearly independent set is linearly independent.\pts{3}
  \item Define Bilinear form and Quadratic form. Let $f: \mathbb{R}^n \times \mathbb{R}^n \to \mathbb{R}$ be defined by $f(x,y) = \alpha_1 \beta_1 + \alpha_2 \beta_2 + \dots + \alpha_n \beta_n$ where $x = (\alpha_1, \dots, \alpha_n)$ and $y = (\beta_1, \dots, \beta_n)$. Show that $f$ is a bilinear form.\pts{3}
\end{parts}

\medskip\rule{\linewidth}{0.2pt}

\noindent\textbf{Question 2.} Matrices \& Normal Form:\pts{12}
\begin{parts}
  \item Show that every square matrix can be expressed uniquely as the sum of a Hermitian matrix and a Skew-Hermitian matrix.\pts{5}
  \item Define rank of a matrix. Reduce the following matrix to normal form and find its rank:\pts{7}
  $$\begin{pmatrix} 1 & 2 & 0 & -1 \\ 3 & 4 & 1 & 2 \\ -2 & 3 & 2 & 5 \end{pmatrix}$$
\end{parts}

\medskip\rule{\linewidth}{0.2pt}

\noindent\textbf{Question 3.} Systems of Linear Equations \& Direct Sums:\pts{12}
\begin{parts}
  \item For what values of $\lambda$, do the equations:
  $$\begin{aligned} x + y + z &= 1 \\ x + 2y + 4z &= \lambda \\ x + 4y + 10z &= \lambda^2 \end{aligned}$$
  have a solution, and solve in each case.\pts{6}
  \item Prove that the vector space $V$ is the direct sum of its subspaces $W_1$ and $W_2$ if and only if: (i) $V = W_1 + W_2$, and (ii) $W_1 \cap W_2 = \{\mathbf{0}\}$.\pts{6}
\end{parts}

\medskip\rule{\linewidth}{0.2pt}

\noindent\textbf{Question 4.} Core Linear Algebra Theorems:\pts{12}
\begin{parts}
  \item State and prove Cayley-Hamilton Theorem.\pts{6}
  \item If $T: V \to U$ is a linear transformation and $V$ is finite dimensional, then prove Rank-Nullity Theorem:\pts{6}
  $$\text{Rank}(T) + \text{Nullity}(T) = \dim(V)$$
\end{parts}

\medskip\rule{\linewidth}{0.2pt}

\noindent\textbf{Question 5.} Change of Basis \& Basis Extension:\pts{12}
\begin{parts}
  \item If the matrix of a linear transformation $T: \mathbb{R}^2 \to \mathbb{R}^2$ with respect to the ordered basis $B = \{(1,0), (0,-1)\}$ is $\begin{pmatrix} 1 & 1 \\ 1 & 1 \end{pmatrix}$, what is the matrix of $T$ with respect to the ordered basis $B' = \{(1,1), (1,-1)\}$?\pts{5}
  \item Let $B = \{(1,1,1,1), (1,2,1,2)\}$ be a linearly independent subset of vector space $\mathbb{R}^4$ over $\mathbb{R}$. Extend it to a basis of $\mathbb{R}^4$.\pts{7}
\end{parts}

\medskip\rule{\linewidth}{0.2pt}

\noindent\textbf{Question 6.} Eigenvalues, Diagonalization \& Gram-Schmidt:\pts{12}
\begin{parts}
  \item For the following matrix, find its eigenvalues and a basis for each of their eigenspaces:
  $$A = \begin{pmatrix} 1 & -3 & 3 \\ 3 & -5 & 3 \\ 6 & -6 & 4 \end{pmatrix}$$
  Is this matrix diagonalizable?\pts{7}
  \item Transform the following basis of Euclidean space $\mathbb{R}^3$ to an orthonormal basis:
  $$\{v_1 = (1,1,1), v_2 = (0,1,1), v_3 = (0,0,1)\}$$
  by applying the Gram-Schmidt orthogonalization process.\pts{5}
\end{parts}

\medskip\rule{\linewidth}{0.2pt}

\noindent\textbf{Question 7.} Inner Product Spaces \& Inequalities:\pts{12}
\begin{parts}
  \item Define unitary linear transformation. Show that $T$ is a unitary linear transformation on an inner-product space $V$ if and only if $\|Tv\| = \|v\|$ for all $v \in V$.\pts{6}
  \item For any two vectors $u, v$ of an inner-product space $V$, state and prove the Cauchy-Schwarz inequality:
  $$|\langle u, v \rangle| \le \|u\| \|v\|$$\pts{6}
\end{parts}

\vfill
\rule{\linewidth}{0.4pt}
\begin{center}\small
\href{https://bhu-exam.vercel.app/}{Click here to visit our website}\\[4pt]
\href{https://me-aryan.vercel.app/}{Click here to visit our contributor}
\end{center}

\end{document}
